\documentclass[10pt,a4paper]{amsart}
\usepackage[margin=19mm]{geometry}
\usepackage{amsmath,amssymb,mathtools}
\usepackage[scr=rsfs,cal=euler]{mathalfa}
\usepackage{xfrac}
\usepackage[unicode,hidelinks,bookmarks=false]{hyperref}
\newcommand{\ie}{i.e.\ }
\newcommand{\cf}{cf.\ }
\newcommand{\resp}{respectively\ }
\newcommand{\Subsection}{\subsection}
\providecommand{\nobreakdash}{\nobreak\hbox{-}}
\setlength{\emergencystretch}{2em}
\multlinegap0pt
\usepackage{amssymb}
\usepackage{mathtools}
\usepackage{tikz}
\usepackage[shortlabels]{enumitem}
\newtheorem{remark}{Remark}
\newtheorem{lemma}{Lemma}
\newtheorem{corollary}{Corollary}
\newtheorem{theorem}{Theorem}
\title{$t$-tone edge coloring of graphs}
\author{Hadeel Al Bazzal}
\date{Source: arXiv:2605.24571v1 (CC BY 4.0)}
\begin{document}
\maketitle

\noindent\emph{Source and licence.} $t$-tone edge coloring of graphs. Version arXiv:2605.24571v1 (CC BY 4.0), distributed by arXiv under the Creative Commons Attribution 4.0 International licence, http://creativecommons.org/licenses/by/4.0/, as recorded in the arXiv licence metadata for this exact version and shown on the abstract page https://arxiv.org/abs/2605.24571v1. Full licence notice: \texttt{licenses/arxiv-2605.24571-CC-BY-4.0.txt}.\par\smallskip

\noindent\textit{Editorial scope.} Counted unit: the article's theorem t5 (source0.tex lines 422--424). The article's complete original proof of it is delivered with it (source0.tex lines 426--455, one \texttt{proof} environment of 30 lines). No mathematics is written, repaired or re-derived: the statement and the proof are the article's own lines, in the article's own notation, and no retained line is substituted or re-flowed.  Retained uncounted as the source-local context the proof needs: lines 195--197; lines 234--236; lines 256--258; lines 265--267; lines 389--391; lines 396--398; lines 407--409.  Nothing was omitted from any retained span, so the package carries no omission record.  No retained line contains a citation, so the wrapper carries no bibliography and no bibliography entry.  No retained line contains a figure environment, an image-inclusion command or drawing code, so no image is delivered and the package declares no asset dependency.  Editorial formatting only, in addition: an AMS \texttt{amsart} preamble that restates the article's own declarations and macros used by the retained lines, namely the environments \texttt{remark}, \texttt{lemma}, \texttt{corollary}, \texttt{theorem} on separate counters, together with the standard packages those lines need.  The retained environments print in this document as Remark 1, Lemma 1, Corollary 1, Theorem 1 in order of appearance, whereas the article prints the same items with the numbering of its own full document.  The complete native source of the article as distributed in the arXiv e-print for this exact version is bundled unchanged as \texttt{sources/201286/source0.tex}.
\begin{remark}\label{r1}
 Let $G'$ be obtained from a graph $G$ by deleting a vertex. Since removing a vertex does not decrease the distance between any pair of edges in $G'$, any $t$-tone edge $k$-coloring of $G'$ is a partial $t$-tone edge $k$-coloring of $G$.
\end{remark}

\begin{remark}\label{r2}
Let $f$ be a partial $2$-tone edge coloring of a graph $G$, and let $e\in E(G)$ be an uncolored edge. If $L$ is a family of pairwise intersecting candidate labels for $e$, then each intermediate vertex of $e$ can be incident to at most one forbidding edge for $e$ under $L$. Otherwise, two such edges would be adjacent, requiring their labels to be disjoint, which contradicts the fact that $L$ is pairwise intersecting.
\end{remark}

\begin{lemma}\label{l1}
Let $f$ be a partial $2$-tone edge coloring of a graph $G$, and let $e$ be an uncolored edge of $G$. Let $L$ be a set of pairwise intersecting candidate labels for $e$. If $|L| > m$, where $m$ is the number of intermediate vertices of $e$, then $f$ can be extended to $e$ in at least $|L| - m$ distinct ways.
\end{lemma}

\begin{corollary}\label{x}
Let $f$ be a partial $2$-tone edge coloring of a graph $G$, and let $e$ be an uncolored edge of $G$. Let $L$ be a set of pairwise intersecting candidate labels for $e$. If $|L| > \eta_G(e)$, then $f$ can be extended to $e$ in at least $|L| - \eta_G(e)$ distinct ways.
\end{corollary}

\begin{lemma}\label{l3}
Let $f$ be a partial $2$-tone edge $7$-coloring of a subcubic graph $G$, and let $e$ be an uncolored edge of $G$. If one end of $e$ has degree one in $G$, then $f$ can be extended to $e$.
\end{lemma}

\begin{lemma}\label{l4}
Let $f$ be a partial $2$-tone edge $10$-coloring of a subcubic graph $G$. If $G$ contains a vertex $u$ of degree two whose incident edges are both uncolored, then $f$ can be extended to these edges.
\end{lemma}

\begin{lemma}\label{l5}
Let $f$ be a partial $2$-tone edge $9$-coloring of a cubic graph $G$. Let $u\in V(G)$ with neighbors $u_1, u_2, u_3$ such that $u_2$ is adjacent to both $u_1$ and $u_3$. If the edges $uu_1, uu_2,$ and $uu_3$ are uncolored, then $f$ can be extended to these edges.
\end{lemma}

\begin{theorem}\label{t5}
Every claw-free subcubic graph admits a $2$-tone edge $11$-coloring.
\end{theorem}

\begin{proof}
Suppose, for a contradiction, that the statement is false, and let $G$ be a counterexample with the minimum number of vertices.

\noindent Our general strategy is to select a vertex $u \in G$ and consider the graph $G' = G - u$. Since the property of being claw-free and subcubic is closed under taking subgraphs, $G'$ is also a claw-free subcubic graph. By the minimality of $G$, the graph $G'$ admits a $2$-tone edge $11$-coloring. By Remark \ref{r1}, this coloring is a partial $2$-tone edge $11$-coloring of $G$ where only the edges incident with $u$ are uncolored. We then demonstrate that this partial coloring can be extended to the uncolored edges, yielding a contradiction. To this end, we first derive several structural properties of $G$.

\noindent Clearly, $G$ is connected; otherwise, by the minimality of $G$, each connected component of $G$ admits a $2$-tone edge $11$-coloring, and the union of these colorings yields a $2$-tone edge $11$-coloring of $G$, a contradiction.

\noindent We next show that $G$ has no vertex of degree one. Suppose, to the contrary, that $u$ is a vertex of degree one in $G$ with neighbor $v$, and let $G' = G - u$. By the minimality of $G$, the graph $G'$ admits a $2$-tone edge $11$-coloring $f$. By Lemma \ref{l3}, $f$ extends to the uncolored edge $uv$, yielding a $2$-tone edge $11$-coloring of $G$, a contradiction.

\noindent Similarly, $G$ has no vertex of degree two. Let $u$ be a vertex of degree two in $G$ with neighbors $v$ and $w$, and let $G' = G - u$. By the minimality of $G$, $G'$ admits a $2$-tone edge $11$-coloring $f$. By Lemma \ref{l4}, $f$ extends to both edges $uv$ and $uw$, again producing a contradiction. Therefore, $G$ is a cubic graph.

\noindent Let $u$ be an arbitrary vertex in $G$ with neighbors $u_1,u_2,u_3$. Since $G$ is claw-free, $\{u_1,u_2,u_3\}$ is not an independent set, and hence at least one pair among these vertices is adjacent. Moreover, $G$ is not isomorphic to $K_4$, since $K_4$ admits a $2$-tone edge $9$-coloring. Therefore, $\{u_1,u_2,u_3\}$ does not induce a triangle in $G$.

\noindent We claim that exactly one pair among $\{u_1,u_2,u_3\}$ is adjacent. Otherwise, suppose that one vertex among $\{u_1,u_2,u_3\}$ is adjacent to the other two; without loss of generality, assume that $u_2$ is adjacent to both $u_1$ and $u_3$. Let $G'=G-u$. By the minimality of $G$, the graph $G'$ admits a $2$-tone edge $11$-coloring $f$. By Lemma~\ref{l5}, $f$ extends to the edges $uu_1$, $uu_2$, and $uu_3$, yielding a contradiction.

\noindent Thus, without loss of generality, we may assume that $u_1u_2\in E(G)$ and that $u_1u_3,u_2u_3\notin E(G)$. Let $G' = G - u$. By the minimality of $G$, $G'$ admits a $2$-tone edge $11$-coloring $f$. Let $u_1'$ and $u_2'$ denote the third neighbors of $u_1$ and $u_2$ in $G$, respectively. We proceed to extend $f$ to $uu_1, uu_2,$ and $uu_3$ in this order.

\noindent Suppose, without loss of generality, that $f(u_1u_2)=12$ and $f(u_1u_1')=34$. Since $d_G(u_1u_1',u_2u_2')\le2$, the label $f(u_2u_2')$ contains a color different from both $3$ and $4$; without loss of generality, assume that it contains $5$. Thus, at least one color from $\{1,2,3,4\}$ is free at $uu_2$; let this color be $3$. Consequently, we assume $f(u_2u_2')$ is either $45$ or $56$.

\noindent It follows that the colors $5,6,7,8,9$ are free at $uu_1$. Consider the family $L=\{56, 57, 58, 59\}$ of pairwise intersecting candidate labels for $uu_1$. Since $\eta_G(uu_1)\le 3$, Corollary~\ref{x} ensures that at least one label in $L$ is valid for $uu_1$. Without loss of generality, we assign $59$ to $uu_1$.

\noindent We next consider the edge $uu_2$. Since the color $5$ appears on both $uu_1$ and $u_2u_2'$, it follows that $f_r(uu_2)\ge 6$. In particular, the colors $3,7,8,10,11$ are free at $uu_2$. Let $L'=\{37,38,310,311\}$ be a family of pairwise intersecting candidate labels for $uu_2$. Let $e$ and $e'$ denote the $uu_3$-edges. By Remark \ref{r2}, at most one of $e$ and $e'$ can be a forbidding edge for $uu_2$ under $L'$.

\noindent Assume, without loss of generality, that $e$ is a forbidding edge for $uu_2$ under $L'$. Since $f(u_1u_1')=34 \notin L'$, and since $uu_2$ has at most two intermediate vertices distinct from $u_1$, Lemma~\ref{l1} implies that at least two labels from $L'$ are valid for $uu_2$. Without loss of generality, assume that these labels are $310$ and $311$. In both labels, $uu_2$ and $e$ share the color $3$. Thus, we have $f_r(uu_3) \ge 4$, yielding at least six candidate labels for $uu_3$.

\noindent Observe that neither $u_1u_1'$ nor $u_2u_2'$ forbids any of these candidate labels for $uu_3$, since $f(u_1u_1')$ contains the color $5$ and $f(u_2u_2')$ contains the color $3$, neither of which appears among the free colors at $uu_3$. Therefore, $uu_3$ has at most five forbidding edges, and hence at least one of its candidate labels is valid. Thus, $f$ extends to the three edges, a contradiction.

\noindent Thus, neither $e$ nor $e'$ forbids a label from $L'$ for $uu_2$. It follows that at least three labels in $L'$ are valid for $uu_2$; assume these are $38, 310,$ and $311$. If any of the colors $3,8,10,11$ appears on $f(e)$ or $f(e')$, then again $f_r(uu_3)\ge 4$, yielding at least six candidate labels for $uu_3$, and as before, $uu_3$ has at most five forbidding edges. Therefore, at least one of these labels is valid, a contradiction. Hence, none of these colors appears on $f(e)$ or $f(e')$.

Now assign $38$ to $uu_2$. Then at least three colors are free at $uu_3$, two of which are $10$ and $11$; let the third be $a$. Let $L'' = \{a10, a11, 1011\}$. Observe that none of the labels of $u_1u_1'$, $u_1u_2$, and $u_2u_2'$ belong to $L''$. Moreover, $uu_3$ has at most two intermediate vertices distinct from $u_1$ and $u_2$. Hence, by Lemma \ref{l1}, at least one label in $L''$ is valid for $uu_3$, and hence $f$ extends to  $uu_1, uu_2,$ and $uu_3$, a contradiction. Thus, such a graph $G$ does not exist.\end{proof}

\end{document}
